% GENERATED FILE. Edit canonical lessons, then run npm run generate:books.
% canonical-source-hash: fnv1a64:346b766bc5a05bf3
% canonical-lessons: ML-C285-paranna, ML-C285-culinna, ML-C285-kuttaneyulla, ML-C285-pollayaya, ML-C285-manniya

\chapter{Flat, Wrinkled, Steep, Hollow, Faded}
\label{ch:ml-flat-wrinkled-steep-hollow-faded}
\hlchaptermodality{285}

\begin{chapteropening}
\textbf{By the end of this chapter:} \emph{I can say flat, wrinkled, steep, hollow and faded.}

Complete the last lesson of chapter 285: I can say flat, wrinkled, steep, hollow and faded.
\end{chapteropening}

\section[\texorpdfstring{\ml{പരന്ന} (paranna)}{പരന്ന (paranna)}]{\ml{പരന്ന} (paranna) — flat}

\label{lesson:ML-C285-paranna}



\begin{quote}
Before the new one: say the Malayalam for boring, then the Malayalam for dangerous.
\end{quote}

\subsection*{You'll want to know: \ml{പരന്ന}}
\textbf{\ml{പരന്ന}} — \emph{paranna} — \textquotedblleft{}flat\textquotedblright{}.

\begin{cousinweb}
\textbf{\ml{പരന്ന}} (\emph{paranna}) is native Dravidian, a past form of the verb \textquotedblleft{}to spread out\textquotedblright{}; DEDR 3949 lists the root with Tamil \emph{para}, Kannada \emph{paraḍu} and Telugu \emph{parapu}.
\end{cousinweb}

\begin{grammarlens}[title={What you've built}]
One more word for this chapter.
\end{grammarlens}

\subsection*{Your turn}
\begin{itemize}
\raggedright
  \item \emph{Say it:} \emph{paranna}
  \item \emph{Say it:} \emph{paranna}, once more
  \item \emph{Say it:} say \emph{paranna}
  \item \emph{Recall:} say the Malayalam for boring, then the Malayalam for dangerous, then say \emph{paranna} again
\end{itemize}

\begin{tcolorbox}[breakable,colback=teal!4,colframe=teal!35!black,title={Before you move on}]
What does \ml{പരന്ന} mean? (\textquotedblleft{}Flat\textquotedblright{}.) Say it once more.
\end{tcolorbox}
\section[\texorpdfstring{\ml{ചുളിഞ്ഞ} (cuḷiñña)}{ചുളിഞ്ഞ (cuḷiñña)}]{\ml{ചുളിഞ്ഞ} (cuḷiñña) — wrinkled, crumpled}

\label{lesson:ML-C285-culinna}



\begin{quote}
Before the new one: say the Malayalam for dangerous, then the Malayalam for flat.
\end{quote}

\subsection*{You'll want to know: \ml{ചുളിഞ്ഞ}}
\textbf{\ml{ചുളിഞ്ഞ}} — \emph{cuḷiñña} — \textquotedblleft{}wrinkled, crumpled\textquotedblright{}.

\begin{cousinweb}
\textbf{\ml{ചുളിഞ്ഞ}} (\emph{cuḷiñña}) is native Dravidian, a past form of the verb \textquotedblleft{}to be wrinkled\textquotedblright{}. DEDR 2700 sets it beside Tamil \emph{cuḷi}, \textquotedblleft{}to be angry, show displeasure\textquotedblright{}, and Tamil \emph{cuḷukku}, \textquotedblleft{}to screw one's face up in displeasure\textquotedblright{}.
\end{cousinweb}

\begin{grammarlens}[title={What you've built}]
One more word for this chapter.
\end{grammarlens}

\subsection*{Your turn}
\begin{itemize}
\raggedright
  \item \emph{Say it:} \emph{cuḷiñña}
  \item \emph{Say it:} \emph{cuḷiñña}, once more
  \item \emph{Say it:} say \emph{cuḷiñña}
  \item \emph{Recall:} say the Malayalam for dangerous, then the Malayalam for flat, then say \emph{cuḷiñña} again
\end{itemize}

\begin{tcolorbox}[breakable,colback=teal!4,colframe=teal!35!black,title={Before you move on}]
What does \ml{ചുളിഞ്ഞ} mean? (\textquotedblleft{}Wrinkled, crumpled\textquotedblright{}.) Say it once more.
\end{tcolorbox}
\section[\texorpdfstring{\ml{കുത്തനെയുള്ള} (kuttaneyuḷḷa)}{കുത്തനെയുള്ള (kuttaneyuḷḷa)}]{\ml{കുത്തനെയുള്ള} (kuttaneyuḷḷa) — steep}

\label{lesson:ML-C285-kuttaneyulla}



\begin{quote}
Before the new one: say the Malayalam for flat, then the Malayalam for wrinkled.
\end{quote}

\subsection*{You'll want to know: \ml{കുത്തനെയുള്ള}}
\textbf{\ml{കുത്തനെയുള്ള}} — \emph{kuttaneyuḷḷa} — \textquotedblleft{}steep\textquotedblright{}.

\begin{cousinweb}
\textbf{\ml{കുത്തനെയുള്ള}} (\emph{kuttaneyuḷḷa}) is \textquotedblleft{}that is steep\textquotedblright{}: \emph{kuttane}, \textquotedblleft{}steeply\textquotedblright{}, a word of uncertain origin, with native \emph{uḷḷa}, \textquotedblleft{}that is\textquotedblright{} (DEDR 697).
\end{cousinweb}

\begin{grammarlens}[title={What you've built}]
One more word for this chapter.
\end{grammarlens}

\subsection*{Your turn}
\begin{itemize}
\raggedright
  \item \emph{Say it:} \emph{kuttaneyuḷḷa}
  \item \emph{Say it:} \emph{kuttaneyuḷḷa}, once more
  \item \emph{Say it:} say \emph{kuttaneyuḷḷa}
  \item \emph{Recall:} say the Malayalam for flat, then the Malayalam for wrinkled, then say \emph{kuttaneyuḷḷa} again
\end{itemize}

\begin{tcolorbox}[breakable,colback=teal!4,colframe=teal!35!black,title={Before you move on}]
What does \ml{കുത്തനെയുള്ള} mean? (\textquotedblleft{}Steep\textquotedblright{}.) Say it once more.
\end{tcolorbox}
\section[\texorpdfstring{\ml{പൊള്ളയായ} (poḷḷayāya)}{പൊള്ളയായ (poḷḷayāya)}]{\ml{പൊള്ളയായ} (poḷḷayāya) — hollow}

\label{lesson:ML-C285-pollayaya}



\begin{quote}
Before the new one: say the Malayalam for wrinkled, then the Malayalam for steep.
\end{quote}

\subsection*{You'll want to know: \ml{പൊള്ളയായ}}
\textbf{\ml{പൊള്ളയായ}} — \emph{poḷḷayāya} — \textquotedblleft{}hollow\textquotedblright{}.

\begin{cousinweb}
\textbf{\ml{പൊള്ളയായ}} (\emph{poḷḷayāya}) is native Dravidian: \emph{poḷḷa}, \textquotedblleft{}a tube; perforated, empty\textquotedblright{} (DEDR 4560), and \emph{āya}, \textquotedblleft{}having become\textquotedblright{} (DEDR 333). DEDR 333 compares Sanskrit \emph{ām}.
\end{cousinweb}

\begin{grammarlens}[title={What you've built}]
One more word for this chapter.
\end{grammarlens}

\subsection*{Your turn}
\begin{itemize}
\raggedright
  \item \emph{Say it:} \emph{poḷḷayāya}
  \item \emph{Say it:} \emph{poḷḷayāya}, once more
  \item \emph{Say it:} say \emph{poḷḷayāya}
  \item \emph{Recall:} say the Malayalam for wrinkled, then the Malayalam for steep, then say \emph{poḷḷayāya} again
\end{itemize}

\begin{tcolorbox}[breakable,colback=teal!4,colframe=teal!35!black,title={Before you move on}]
What does \ml{പൊള്ളയായ} mean? (\textquotedblleft{}Hollow\textquotedblright{}.) Say it once more.
\end{tcolorbox}
\section[\texorpdfstring{\ml{മങ്ങിയ} (maṅṅiya)}{മങ്ങിയ (maṅṅiya)}]{\ml{മങ്ങിയ} (maṅṅiya) — faded, dim}

\label{lesson:ML-C285-manniya}



\begin{quote}
Before the new one: say the Malayalam for steep, then the Malayalam for hollow.
\end{quote}

\subsection*{You'll want to know: \ml{മങ്ങിയ}}
\textbf{\ml{മങ്ങിയ}} — \emph{maṅṅiya} — \textquotedblleft{}faded, dim\textquotedblright{}.

\begin{cousinweb}
\textbf{\ml{മങ്ങിയ}} (\emph{maṅṅiya}) is native Dravidian, a past form of the verb \textquotedblleft{}to grow dim, pale\textquotedblright{}; DEDR 4750 lists the root with Tamil \emph{maḻu} and Kannada \emph{maḻa}. DEDR 4750 compares Sanskrit \emph{maṅku}.
\end{cousinweb}

\begin{grammarlens}[title={What you've built}]
One more word for this chapter.
\end{grammarlens}

\subsection*{Your turn}
\begin{itemize}
\raggedright
  \item \emph{Say it:} \emph{maṅṅiya}
  \item \emph{Say it:} \emph{maṅṅiya}, once more
  \item \emph{Say it:} say \emph{maṅṅiya}
  \item \emph{Recall:} say the Malayalam for steep, then the Malayalam for hollow, then say \emph{maṅṅiya} again
\end{itemize}

\begin{tcolorbox}[breakable,colback=teal!4,colframe=teal!35!black,title={Before you move on}]
What does \ml{മങ്ങിയ} mean? (\textquotedblleft{}Faded, dim\textquotedblright{}.) Say it once more.
\end{tcolorbox}
